\clearpage
\VerdierPhysicalPageAnchor{0249}%
\begingroup\small\setlength{\parskip}{0pt}\setlength{\emergencystretch}{1.5em}
\setlength{\abovedisplayskip}{4pt plus 2pt minus 2pt}%
\setlength{\belowdisplayskip}{4pt plus 2pt minus 2pt}%

\VerdierSemanticAnchor{ch4d-4-6-9-first-differentials}
\noindent 而第一微分为 (4.1.9.1)
\[
\begin{aligned}
d_1^{p,q}\colon(\mathrm{C.E.})_1^{p,q}
  &\longrightarrow(\mathrm{C.E.})_1^{p+1,q},\\
d_1^{p,q}&=F^q(-d_2^{\bullet,p}).
\end{aligned}
\]

\VerdierSemanticAnchor{ch4d-4-6-9-column-decomposition}
复形 $I^{\bullet,p}$ 是 Cartan--Eilenberg 内射复形 (4.6.1)，因此由 (4.6.3)，
$D(\mathcal A)$ 中存在典范同构
\[
I^{\bullet,p}\xrightarrow{\sim}
\prod_{r\in\mathbb Z}H^r(I^{\bullet,p})[-r].
\]

\VerdierSemanticAnchor{ch4d-4-6-9-h1-components}
现在采用 (4.6.4.3) 中引入的复形 $H_1^r(I^{\bullet\bullet})$。
对每个整数 $r$，都有
\[
H_1^r(I^{\bullet\bullet})^p=H^r(I^{\bullet,p}),
\]
因而 $D(\mathcal A)$ 中存在典范同构

\VerdierSemanticAnchor{ch4d-4-6-9-h1-decomposition}
\[
I^{\bullet,p}\xrightarrow{\sim}
\prod_{r\in\mathbb Z}H_1^r(I^{\bullet\bullet})^p[-r],
\]
并且，在这个同构之下，复形态射
\[
d_2^{\bullet,p}\colon
I^{\bullet,p}\longrightarrow I^{\bullet,p+1}
\]
同构于态射

\VerdierSemanticAnchor{ch4d-4-6-9-d2-decomposition}
\[
\prod_{r\in\mathbb Z}
d_{H_1^r(I^{\bullet\bullet})}^{p}[-r],
\]
由此得到初始项的最终表达式

\VerdierSemanticAnchor{ch4d-4-6-9-1}
\[
(\mathrm{C.E.})_1^{p,q}
=F^q\left(
\prod_{r\in\mathbb Z}H_1^r(I^{\bullet\bullet})^p[-r]
\right),
\tag{4.6.9.1}
\]
以及第一微分的最终表达式

\VerdierSemanticAnchor{ch4d-4-6-9-2}
\[
d_1^{p,q}
=-F^q\left(
\prod_{r\in\mathbb Z}d_{H_1^r(I^{\bullet\bullet})}^{p}[-r]
\right).
\tag{4.6.9.2}
\]

\VerdierSemanticAnchor{ch4d-4-6-9-3-prelude}
设 $F^{q,\bullet}(H_1^r(I^{\bullet\bullet}))$ 是把函子 $F^q$ 施于复形
$H_1^r(I^{\bullet\bullet})$ 的各分量和微分而得到的 $\mathcal B$ 中对象的
复形。公式 (4.6.9.1) 与 (4.6.9.2) 定义了一个态射

\VerdierSemanticAnchor{ch4d-4-6-9-3}
\[
(\mathrm{C.E.})_2^{p,q}\longrightarrow
\prod_{r\in\mathbb Z}
H_{\mathcal B}^{p}\left(
F^{q-r,\bullet}(H_1^r(I^{\bullet\bullet}))
\right).
\tag{4.6.9.3}
\]
\endgroup
